\clearpage
\section{Ordered distributors and graph assembly}
\label{sec:copying-assembly}
The next prerequisite is a concrete construction of the distributors of
Definition~10.1. We assemble actual finite typed graphs, with the paper's
ordered interfaces, and prove that assembly preserves incidence and paths.
These results provide wiring for the future translation; they do not yet
supply copying transitions, store evaluation, template robustness, or
Proposition~12.2.

\subsection{Parallel assembly with an ordered boundary}
\termdef{ordered-parallel-assembly}
Let $A$ be a finite index set and let $G_a:I_a\to O_a$ be typed graphs.
A boundary specification consists of type-preserving bijections from the
disjoint unions of their input and output ports to the chosen ordered
interfaces $J$ and $K$. Thus every component boundary port occurs exactly once.
This supports interleaving, rather than merely concatenating interfaces.

The boundary structure records these bijections and their typing equations.
The assembled vertex type consists of tagged internal vertices and the new
external ports. The embedding tags an internal vertex by $a$ and sends an
external port through the appropriate boundary bijection. Its inverse
recovers the tag and local port. The inverse equations follow by cases on
internal, input, and output vertices and by the two boundary inverse laws.

\lean{Graph/Parallel}{15}{21}{Parallel.Boundary}

\lean{Graph/Parallel}{23}{23}{Parallel.Vertex}

\lean{Graph/Parallel}{26}{29}{Parallel.embed}

\lean{Graph/Parallel}{32}{51}{Parallel.vertices}

Embedding preserves whether a vertex is an input or output, as a direct
case analysis shows. It is injective because the joint vertex map is a
bijection. It preserves wire labels: internal labels are inherited, and
external labels are preserved by the boundary specification.

\lean{Graph/Parallel}{53}{54}{Parallel.embed_isInput}

\lean{Graph/Parallel}{56}{57}{Parallel.embed_isOutput}

\lean{Graph/Parallel}{59}{63}{Parallel.embed_injective}

\lean{Graph/Parallel}{65}{71}{Parallel.embed_label}

For an edge $(a,e)$, transport its incidence list through the component
embedding. Injectivity preserves absence of repetitions. If two transported
incidence lists contain the same vertex, the joint inverse first identifies
the component tags and then the original vertex; local uniqueness therefore
identifies the edges. For coverage, a local witness transports forward.
Conversely, take the unique preimage of a global vertex and use local
coverage there. This works for either source or target incidence by
parameterising the vertex predicate.

\lean{Graph/Parallel}{74}{75}{Parallel.incidence}

\lean{Graph/Parallel}{77}{79}{Parallel.incidence_nodup}

\lean{Graph/Parallel}{81}{95}{Parallel.incidence_unique}

\lean{Graph/Parallel}{98}{114}{Parallel.incidence_coverage}

The assembled graph has edges and internal vertices given by dependent
disjoint unions. Finite enumerations are formed from the enumeration of $A$
and the component enumerations. Source and target incidence use the preceding
transport. Their uniqueness, coverage, and absence of duplicates are exactly
the lemmas just proved. Typing follows by mapping the label-preservation
equation over each incidence list. Thus this construction satisfies all
finite-graph fields, including the exhaustive boundary laws.

\lean{Graph/Parallel}{117}{145}{Parallel.net}

If two assembled edges are adjacent, a shared target/source vertex has
preimages in both components. Injectivity forces the tags to agree.
Within a component, adjacency is equivalent to the original adjacency:
reflect the shared vertex using injectivity, or transport it forward.
Induction on a nonempty path proves that all its edges have the same tag.
A second induction proves preservation and reflection of paths under a
component inclusion. Empty lists are excluded by the existing path definition;
single-edge paths give the base case. Finally a label predicate holds of
all assembled edges precisely when it holds of every component edge,
since edges are tagged pairs and their labels are unchanged.

\lean{Graph/Parallel}{147}{153}{Parallel.adjacent_tag}

\lean{Graph/Parallel}{156}{166}{Parallel.adjacent_same}

\lean{Graph/Parallel}{169}{188}{Parallel.path_tags}

\lean{Graph/Parallel}{191}{205}{Parallel.path_same}

\lean{Graph/Parallel}{208}{212}{Parallel.labels_iff}

\subsection{The finite input permutation and distributor forest}
\termdef{copying-distributor}
For $k$ resources each used $m$ times, the distributor $D_{k,m}$ consists of
$k$ canonical trees. In zero-based indexing, resource $i$ at use $j$ has input
position $i+kj$ and output position $i$. This is the ordered interface of
Definition~10.1, not just an equality of interface sizes.

The grid index maps $(i,j)$ to $i+kj$. Its bound follows from $i<k$ and
$j<m$. The inverse sends $n$ to $(n\bmod k,\lfloor n/k\rfloor)$.
The existence of $n<km$ implies $k>0$, so the division bounds are valid
even though the construction allows an empty grid. The inverse equations
use Euclidean division, $i\bmod k=i$, and $\lfloor i/k\rfloor=0$.
They give a bijection. The final remainder and quotient lemmas explicitly
recover the resource and use coordinates.

\lean{Graph/FinGrid}{8}{12}{FinGrid.index}

\lean{Graph/FinGrid}{15}{21}{FinGrid.split}

\lean{Graph/FinGrid}{23}{29}{FinGrid.split_index}

\lean{Graph/FinGrid}{31}{33}{FinGrid.index_split}

\lean{Graph/FinGrid}{35}{39}{FinGrid.bijection}

\lean{Graph/FinGrid}{42}{43}{FinGrid.index_mod}

\lean{Graph/FinGrid}{46}{47}{FinGrid.index_div}

The distributor input functions are this grid bijection with the list-length
casts needed by typed interfaces. Their inverse proofs reduce to the same
remainder and quotient equations. Each tree has one output, so the output
bijection forgets its unique local output index. Both boundary label equations
hold because all involved ports have the chosen sharing sort.

\lean{Copying/Distributor}{14}{18}{Distributor.inputIndex}

\lean{Copying/Distributor}{20}{23}{Distributor.inputSplit}

\lean{Copying/Distributor}{25}{35}{Distributor.inputSplit_index}

\lean{Copying/Distributor}{37}{40}{Distributor.inputIndex_split}

\lean{Copying/Distributor}{43}{58}{Distributor.boundary}

Instantiate parallel assembly with $k$ copies of the canonical $m$-input
tree. The input- and output-position equations follow by unfolding the
boundary maps. An edge $(i,e)$ is a weakening exactly when $e=0$, and is
a contraction exactly when $e\ne0$, by the canonical-tree label lemmas.
Relabelling primitive labels embeds this graph into every finite box allowance;
it does not create any boxes or operational transitions.

\lean{Copying/Distributor}{61}{63}{Distributor.net}

\lean{Copying/Distributor}{66}{68}{Distributor.input_position}

\lean{Copying/Distributor}{70}{71}{Distributor.output_position}

\lean{Copying/Distributor}{74}{76}{Distributor.weakening_iff}

\lean{Copying/Distributor}{78}{80}{Distributor.contraction_iff}

\lean{Copying/Distributor}{83}{86}{Distributor.atDepth}

A general rank argument supplies acyclicity. Suppose a natural-number rank
strictly increases across adjacency. Induction on a path, using transitivity
at its head, proves that the rank at any earlier position is less than the
rank at any later position. If an edge occurred twice, compare its two
positions; either strict ordering would make its rank less than itself.
Hence every path has no repeated edge.

\lean{Graph/PathRank}{11}{37}{PathRank.path_rank_strict}

\lean{Graph/PathRank}{39}{50}{PathRank.path_nodup_of_rank}

The generic parallel lemmas show that adjacent distributor edges belong
to the same tree and that their local edges are adjacent there. Therefore
the canonical-tree rank remains strictly increasing on distributor adjacency.
Apply the general rank theorem to obtain absence of repeated edges on every
path. Interleaving boundary ports has introduced no path between resources.

\lean{Copying/Distributor}{89}{91}{Distributor.adjacent_tag}

\lean{Copying/Distributor}{93}{95}{Distributor.adjacent_same}

\lean{Copying/Distributor}{98}{105}{Distributor.adjacent_rank_lt}

\lean{Copying/Distributor}{108}{111}{Distributor.path_nodup}

When $k=0$, there is no edge or internal vertex: either would require an
element of the empty resource index set. When $m=0$, each component has only
its edge numbered zero, so every edge is a weakening. Thus zero uses does
not silently discard resources. More generally, any primitive-label property
true of weakening and contraction holds throughout the depth-indexed
distributor. Unfold its labels and split on whether the local edge number
is zero. Since neither constructor is a focus or hole, this also proves
global focus freedom and hole freedom for both term and atom distributors.

\lean{Copying/Distributor}{114}{114}{Distributor.zero_edges}

\lean{Copying/Distributor}{116}{116}{Distributor.zero_internal}

\lean{Copying/Distributor}{119}{125}{Distributor.zero_uses}

\lean{Copying/Distributor}{128}{133}{Distributor.allLabels_atDepth}

\lean{Copying/Distributor}{135}{136}{Distributor.focusFree}

\lean{Copying/Distributor}{138}{139}{Distributor.holeFree}

\subsection{Serial composition through a shared interface}
\termdef{copying-serial-assembly}
Given $G:I\to J$ and $H:J\to K$, serial composition uses a two-hole pattern
whose internal vertices are exactly the ordered ports of $J$. The left hole
has interface $I\to J$; the right hole has interface $J\to K$. Their hole
identities differ even when their interfaces coincide.

The pattern has two edges. Its left edge runs from outer inputs to shared
internal vertices; its right edge runs from those internal vertices to outer
outputs. Enumerating the corresponding finite index lists establishes all
incidence and typing laws. Distinct port constructors exclude false matches
between outer and internal vertices. For every $j$ in $J$, the left target
port and right source port are the same internal vertex; the connection
lemma proves this by evaluating the indexed incidence lists. The construction
also handles an empty shared interface.

\lean{Copying/Serial}{11}{13}{Serial.hole}

\lean{Copying/Serial}{16}{80}{Serial.pattern}

\lean{Copying/Serial}{83}{86}{Serial.interface_connection}

Fill the left slot with $G$ and the right with $H$, using the signature
equalities to check their interfaces. Simultaneous expansion replaces both
pattern edges and glues their boundary vertices at the pattern ports.
This defines the composed graph; the pattern holes are no longer edges.
For a primitive-label predicate, each expanded edge comes from one filling,
so its label property is inherited from that filling. Specialise this fact
to the two slots, then to focus freedom and hole freedom. These are structural
closure laws, not a theorem that an operational reduction can cross an
arbitrary surrounding context.

\lean{Copying/Serial}{89}{94}{Serial.filling}

\lean{Copying/Serial}{97}{98}{Serial.net}

\lean{Copying/Serial}{101}{106}{Serial.allLabels_expand}

\lean{Copying/Serial}{109}{113}{Serial.allLabels_net}

\lean{Copying/Serial}{115}{116}{Serial.focusFree_net}

\lean{Copying/Serial}{118}{119}{Serial.holeFree_net}

\subsection{Concrete checks and remaining obligations}

The input-order regression computes all six positions for three resources
used twice: resource zero occupies positions $0,3$, resource one $1,4$,
and resource two $2,5$. The separation check rules out adjacency between
the two distinct resources in a two-resource distributor by tag equality.
A single-use tree still contains its weakening and contraction; it is not
replaced by an identity. The zero-use check retains two distinct weakenings
for two resources. These label calculations hold for both sharing sorts.
Finally, serially assemble an instance edge from term to atom with a store
edge from atom to term. The resulting term-to-term graph is focus-free and
hole-free by the serial closure lemmas and inspection of the two labels.
This last example checks atom-interface wiring, not allocation, reading,
or writing in a running machine.

\lean{Copying/AssemblyExamples}{8}{15}{AssemblyExamples.distributor_input_order}

\lean{Copying/AssemblyExamples}{18}{22}{AssemblyExamples.distributor_separation}

\lean{Copying/AssemblyExamples}{25}{28}{AssemblyExamples.distributor_single_use}

\lean{Copying/AssemblyExamples}{31}{34}{AssemblyExamples.distributor_no_uses}

\lean{Copying/AssemblyExamples}{37}{38}{AssemblyExamples.instanceStore}

\lean{Copying/AssemblyExamples}{40}{42}{AssemblyExamples.instanceStore_focusFree}

\lean{Copying/AssemblyExamples}{44}{46}{AssemblyExamples.instanceStore_holeFree}

The full build and all three axiom audits pass. The new audited declarations
use only Lean's standard propositional extensionality, classical choice, and
quotient soundness. No machine determinism, template robustness, or
translation-correctness assumption has been introduced. Arbitrary contraction
trees are now treated in Section~\ref{sec:copying-tracing}. Copyability and stability are treated in Section~\ref{sec:copying-values}.
Context translation and the translation laws,
and the operational and template arguments remain to be proved.
